\documentclass[11pt]{article}
\usepackage[margin=1in]{geometry}
\usepackage{amsmath,amssymb}
\usepackage[hidelinks]{hyperref}

\title{Literature Status: The Garling Symmetric-Basis Question}
\author{Run fa\_banach\_001}
\date{2026-06-28}

\begin{document}
\maketitle

\paragraph{Status.}
\texttt{literature\_already\_answered}.  This packet records a separate later
paper that explicitly answers the extracted question from arXiv:1703.07772.

\paragraph{Original source and question.}
In Section 7 of Fernando Albiac, Jose L. Ansorena, and Ben Wallis,
\emph{On Garling sequence spaces}, arXiv:1703.07772, the authors close with
the question:
\[
\text{Do there exist } \mathbf{w}\in\mathcal W \text{ and } 1\le p<\infty
\text{ so that } g(\mathbf{w},p) \text{ admits a symmetric basis?}
\]

\paragraph{Answering source.}
Fernando Albiac, Jose L. Ansorena, Denny Leung, and Ben Wallis,
\emph{Optimality of the rearrangement inequality with applications to
Lorentz-type sequence spaces}, arXiv:1705.03936, explicitly identifies this
question as answered.  Its abstract says that, as a byproduct, the paper
settles a problem raised in the 1703.07772 paper and proves that Garling
sequence spaces have no symmetric basis.  The introduction repeats that
its main rearrangement theorem is applied to show that Garling sequence spaces
have no symmetric basis, answering the problem posed in the earlier paper.

\paragraph{Exact answer.}
Section 3 of arXiv:1705.03936 is titled ``Applications to Garling sequence
spaces.''  Lemma 3.1 states that the canonical sequence is not a symmetric
basis for $g(\mathbf{w},p)$ for any $\mathbf{w}\in\mathcal W$ and any
$1\le p<\infty$.  Theorem 3.2(i) states the full structural conclusion:
\[
\text{There is no symmetric basis for } g(\mathbf{w},p).
\]
Thus the existence question from arXiv:1703.07772 has a complete negative
answer.

\paragraph{Scope.}
The supporting authors explicitly knew they were answering the source question.
This packet makes no new mathematical claim beyond recording that the target is
already settled in the literature.  No part of the extracted existence question
remains open here.

\begin{thebibliography}{9}
\bibitem{AAW}
F. Albiac, J. L. Ansorena, and B. Wallis,
\emph{On Garling sequence spaces}, arXiv:1703.07772.

\bibitem{AALW}
F. Albiac, J. L. Ansorena, D. Leung, and B. Wallis,
\emph{Optimality of the rearrangement inequality with applications to
Lorentz-type sequence spaces}, arXiv:1705.03936.
\end{thebibliography}

\end{document}
